\documentclass{article}
\usepackage{amsmath}
\usepackage{amssymb}
\usepackage{hyperref}

\title{Review of "Measuring Sample Quality with Stein’s Method"}
\author{Elham Afzali}

\begin{document}
	
	\maketitle
	
	\section*{Paper Information}
	\textbf{Title:} Measuring Sample Quality with Stein’s Method \\
	\textbf{Authors:} Jackson Gorham, Department of Statistics, Stanford University \\
	\textbf{Authors:} Lester Mackey, Department of Statistics, Stanford University \\
	\textbf{Published in:} [Journal/Conference Name, Year (if available)] \\
	\textbf{Link:} [Provide link if available]
	
	\section*{Introduction}
	The paper "Measuring Sample Quality with Stein’s Method" by Jackson Gorham and Lester Mackey introduces an innovative approach to assess the quality of samples generated by Markov Chain Monte Carlo (MCMC) methods, particularly those that intentionally incorporate asymptotic bias to enhance computational efficiency. This paper addresses a critical gap in traditional sample quality measures, such as effective sample size, which do not account for the asymptotic bias introduced by certain MCMC procedures.
	
	\section*{Background}
	Markov Chain Monte Carlo (MCMC) methods are commonly employed to approximate intractable expectations of the form:
	
	\[
	\mathbb{E}_P[h(Z)] = \int_X p(x)h(x)dx
	\]
	
	where \(P\) is the target distribution, and \(h\) is a test function. In practice, we approximate \(\mathbb{E}_P[h(Z)]\) using samples from an MCMC algorithm, leading to estimates of the form:
	
	\[
	\mathbb{E}_Q[h(X)] = \frac{1}{n} \sum_{i=1}^n h(X_i)
	\]
	
	where \(Q\) denotes the empirical distribution of the samples. However, traditional MCMC methods assume eventual convergence to the target distribution \(P\). The introduction of asymptotic bias in recent MCMC methods has posed challenges for these conventional sample quality measures.
	
	\section*{Contributions}
	To address these challenges, the authors propose a new sample quality measure based on Stein's method. This measure is designed to quantify the discrepancy between the empirical distribution \(Q\) and the target distribution \(P\) by evaluating the maximum deviation between their expectations over a large class of test functions \(H\):
	
	\[
	d_H(Q, P) = \sup_{h \in H} \left| \mathbb{E}_Q[h(X)] - \mathbb{E}_P[h(Z)] \right|
	\]
	
	The innovation in this paper is the use of Stein's method to define a discrepancy measure, known as the Stein discrepancy:
	
	\[
	S(Q, \mathcal{T}, \mathcal{G}) = \sup_{g \in \mathcal{G}} \left| \mathbb{E}_Q[\mathcal{T}g(X)] - \mathbb{E}_P[\mathcal{T}g(Z)] \right|
	\]
	
	where \(\mathcal{T}\) is a Stein operator associated with the target distribution \(P\), and \(\mathcal{G}\) is a suitable class of test functions. This measure has the advantage of being computable without requiring explicit integration under \(P\).
	
	\section*{Strengths}
	\begin{itemize}
		\item \textbf{Theoretical Rigor:} The authors provide a strong theoretical foundation for their measure, showing that the Stein discrepancy can be used to assess the convergence of MCMC samples even when traditional measures fail. They relate the Stein discrepancy to well-known probability metrics like the Wasserstein distance, further solidifying its theoretical underpinnings.
		\item \textbf{Computational Efficiency:} The paper introduces a streamlined computational approach using geometric spanners, which allows for efficient computation of the Stein discrepancy in high-dimensional settings.
		\item \textbf{Applicability:} The authors demonstrate the utility of the Stein discrepancy in practical applications, such as hyperparameter selection and convergence rate assessment, making it a valuable tool for practitioners.
	\end{itemize}
	
	\section*{Weaknesses}
	\begin{itemize}
		\item \textbf{Implementation Complexity:} Despite the streamlined approach, implementing the Stein discrepancy measure may still be challenging for practitioners who are not familiar with Stein’s method or the associated computational techniques like geometric spanners.
		\item \textbf{Example Diversity:} The empirical evaluation is robust but could be further strengthened by including a more diverse set of examples to showcase the versatility of the Stein discrepancy in various sampling scenarios.
	\end{itemize}
	
	\section*{Conclusion}
	"Measuring Sample Quality with Stein’s Method" makes a significant contribution to the field of MCMC sampling by introducing a robust and theoretically sound measure for evaluating sample quality. The Stein discrepancy provides a powerful tool for comparing exact, biased, and deterministic sample sequences, addressing a crucial need in modern computational statistics. While the method may present implementation challenges, its potential impact on the field is undeniable.
	
\end{document}
